% appendix.tex — technical appendix: preliminaries and all proofs
% Source: docs/math_appendix_v1_en.md (authoritative English edition; faithful; skeleton-status items marked per source).
\section{Preliminaries: directed information and testing duality}
\label{app:prelim}

We write $X^T=(X_1,\dots,X_T)$; all logarithms are natural (nats). Standard
facts (chain rule, DPI, concavity of mutual information in the input
distribution, strict concavity of entropy, Pinsker's inequality) are used as
in \citet{cover_thomas2006}.

\begin{definition}[Directed information; Massey 1990, eq.~(7) \citep{massey1990}]
$I(X^T\!\to\!Y^T):=\sum_{t=1}^T I(X^t;Y_t\given Y^{t-1})$.
\end{definition}

\begin{definition}[Causally conditioned directed information; Kramer 2003,
\S II-C Def.~4, eq.~(4) \citep{kramer2003}]
$I(X^T\!\to\!Y^T\,\|\,Z^T):=\sum_{t=1}^T I(X^t;Y_t\given Y^{t-1},Z^t)$.
This \emph{non-anticipating} conditioning ($Z^t$) is the convention used
throughout the paper: it is what makes the mediation lemma
(Lemma~\ref{lem:mediation}) true, since conditioning on future components of
$Z$ can open collider paths through their descendants
(Appendix~\ref{app:t2}).\footnote{We verified against the published original
that Kramer's Definition~4 is the causal variant with $Z^t$ conditioning; the
full-sequence variant is Marko's convention. A direction misprint in Kramer's
printed Property~2 counterexample was verified and corrected: with $X_1\perp
Y_1$ fair bits, $X_2=X_1$, $Y_2=X_1\oplus Y_1$, one has
$I(Y^2\!\to\!X^2)=0$ while $I(X^2;Y^2)=1$ bit---the printed sentence reverses
the direction of the first DI; our regression tests encode the corrected
version.}
\end{definition}

\begin{lemma}[MI = DI under randomized exogenous probes]\label{lem:probe}
Let $I^T$ be exogenous to the agent (experimenter-randomized,
$I^T\ind$ agent history) and non-anticipative (given $I^t$, the choice of
$I_{t+1:T}$ is independent of $A^t$). Then
$I(I^T;A^T)=I(I^T\!\to\!A^T)$.
\end{lemma}
\begin{proof}
$\ge$: $I(I^T;A^T)=\sum_t I(I^T;A_t\given A^{t-1})\ge\sum_t
I(I^t;A_t\given A^{t-1})=I(I^T\!\to\!A^T)$ since conditioning reduces MI and
$I^t\subseteq I^T$. $\le$: it suffices that
$I(I_{t+1:T};A_t\given I^t,A^{t-1})=0$: causality makes $A_t$'s generation
independent of $I_{t+1:T}$, and randomization makes $I_{t+1:T}$ independent of
$(A^t,A_t)$ given $I^t$. The $\le$ direction is the randomized-probe instance
of Massey's no-feedback identity (mutual information equals directed
information when the input process carries no feedback from the output):
randomization supplies the no-feedback premise that the closed agent loop
would otherwise violate.
\end{proof}

\begin{remark}[From the probe lemma to the certified bound: the bypass
premise]\label{rem:probe-bridge}
Lemma~\ref{lem:probe} identifies the observable $I(I^T;A^T)$ with the
\emph{unconditioned} directed information $I(I^T\!\to\!A^T)$; the estimand of
the certified bound is the \emph{causally conditioned}
$\Fext=I(I^T\!\to\!A^T\|S^T)$. The bridge is the \emph{bypass premise}
$I^t\ind S^t\given A^{t-1}$ (the probe acts outside the belief-mediation
path). Under it, $I(I^t;S^t\given A^{t-1})=0$, so the identity
\[
I(I^t;A_t\given A^{t-1},S^t)-I(I^t;A_t\given A^{t-1})
= I(I^t;S^t\given A^{t-1},A_t)-I(I^t;S^t\given A^{t-1})
\]
has a non-negative right-hand side for every $t$: conditioning on $S^t$
is non-decreasing, and $F_{ext}^{obs}\le\Fext$. Without the premise the
direction can reverse; the mediation witness ($S=I\oplus N$ with
$N\sim\mathrm{Bern}(0.1)$, $A=S$: $\Fext=0$ while $I(I;A)=0.531$~bit) has the
probe entering through $S$, and then the plug-in overestimates the external
flow. The premise is not a regularity condition but the \emph{semantic
certification boundary} of \eqref{eq:aorg-interval}: only bypass-semantics
probes carry the one-sided reading, and adaptive (state-correlated) probes
violate it (BB-2). We register the boundary rather than extend the theorem
over the mediation case, where no black-box direction is available.
\end{remark}

\begin{lemma}[Testing duality]\label{lem:tv-testing}
For any test $\psi\in[0,1]$ (randomization allowed):
$P_0(\psi=1)+P_\gamma(\psi=0)\ge 1-\TV(P_0,P_\gamma)$.
\end{lemma}
\begin{proof}
$P_0(\psi=1)-P_\gamma(\psi=1)\le\TV$, so
$P_0(\psi=1)+1-P_\gamma(\psi=1)\ge 1-\TV$.
\end{proof}

\begin{corollary}[Le Cam-type undetectability]
If $\TV(P_\gamma,P_0)\le\varepsilon$, every level-$\alpha$ test has power
$\le\alpha+\varepsilon$.
\end{corollary}

\begin{lemma}[Perturbation bound for contraction fixed points]
\label{lem:contraction}
Let $T_\theta:\Theta\to\Theta$ be a contraction with modulus $c<1$ in a common
metric $d$, with $d(T_\theta(x),T_{\theta'}(x))\le L\,d_\Theta(\theta,\theta')$.
Then the unique fixed points satisfy
$d(u^*(\theta),u^*(\theta'))\le\frac{L}{1-c}d_\Theta(\theta,\theta')$.
\end{lemma}
\begin{proof}
$d(u^*,u'^*)\le d(u^*,T_\theta(u'^*))+d(T_\theta(u'^*),T_\theta(u^*))
+d(T_\theta(u^*),u'^*)\le L\,d_\Theta(\theta,\theta')+c\,d(u^*,u'^*)$;
rearrange. This lemma bounds goal drift by belief drift (the $\gamma\to\Agoal$
transmission) via $T=\Endorse_{B,M}$.
\end{proof}

\section{The representation theorem: full proofs}\label{app:representation}

Recall (A1)--(A6) from Definition~\ref{def:a16}.

\begin{lemma}[Quasilinear representation, cited]\label{lem:quasilinear}
\textup{(A1)}$+$\textup{(A2)}$+$\textup{(A3)} imply the existence of a
continuous strictly increasing $\phi:(0,1]\to\mathbb{R}$ and $w_i>0$
($\sum_iw_i=1$) with $F(x)=\phi^{-1}\bigl(\sum_iw_i\phi(x_i)\bigr)$.
\begin{proof}[How the citation is used]
\textup{(A3)} supplies, for \emph{every} non-trivial partition $\{K,L\}$,
continuous increasing $F_K,F_L,G_{KL}$ with $F=G_{KL}(F_K(x_K),F_L(x_L))$; the
cross-partition consistency of this family (the same $F$ through every
grouping) is precisely the hypothesis of the overlapping-decomposition and
separability theorems, which reduce coordinate-dependent scales to a single
generator: Debreu's separability representation gives coordinate scales, and
the reduction to a common generator with constant weights is the
Acz\'el--Gorman step (\citealp{debreu1960}; \citealp{aczel1966}, Ch.~6;
\citealp{aczel_dhombres1989}, Ch.~17; \citealp{gorman1968}; for $n=2$ the
classical \citealp{nagumo1930,kolmogorov1930} result). The quantifier
alignment matters and is satisfied by construction: \textup{(A3)} quantifies
over all partitions of one fixed $F$, and the cited theorems return one
generator $\phi$ common to all subgroups. Biseparability is exactly the
needed weakening of associativity.
\end{proof}
\end{lemma}

\begin{lemma}[AX5 $\iff$ $\phi(0^+)=-\infty$]\label{lem:phi0}
\end{lemma}
\begin{proof}
($\Rightarrow$) Take $x_1\to0^+$ with other coordinates fixed: $F\to0$ by (A4)
and continuity. Since $F=\phi^{-1}(w_1\phi(x_1)+C)$ with $C$ finite, if
$\phi(0^+)>-\infty$ then $F\to\phi^{-1}(w_1\phi(0^+)+C)>0$, contradiction.
($\Leftarrow$) A coordinate $\to0$ sends the sum to $-\infty$, whence
$F\to\phi^{-1}(-infty):=0$; $\min>0$ keeps the sum finite, whence $F>0$.
\end{proof}

\begin{lemma}[Homogeneity $\Rightarrow$ power family]\label{lem:power}
Under Lemmas~\ref{lem:quasilinear}--\ref{lem:phi0} plus \textup{(A5)}, there
exists $p\le0$ with $F=M_p$ ($p=0$ the geometric mean; for $p<0$, $\phi$
affinely equivalent to $(x^p-1)/p$).
\end{lemma}
\begin{proof}
(1) Write $J:=\phi((0,1])$ (an interval with lower endpoint $-\infty$) and
$h_t(u):=\phi(t\phi^{-1}(u))$. Applying $\phi$ to $F(tx)=tF(x)$ gives
$\sum_iw_ih_t(\phi(x_i))=h_t(\sum_iw_i\phi(x_i))$: $h_t$ satisfies the
weighted Jensen equation on $J$; continuity makes $h_t$ affine,
$h_t(u)=a(t)u+b(t)$, $a(t)>0$.
(2) Back in $x$: $\phi(tx)=a(t)\phi(x)+b(t)$. Set
$\Phi(x):=\phi(x)-\phi(1)$ and $B(t):=b(t)+(a(t)-1)\phi(1)$; evaluating at
$x=1$ gives $B(t)=\Phi(t)$, hence ($\star$) $\Phi(tx)=A(t)\Phi(x)+\Phi(t)$.
(3) Symmetry $\Phi(tx)=\Phi(xt)$ with ($\star$) gives
$\Phi(x)(A(t)-1)=\Phi(t)(A(x)-1)$; for $x,t\neq1$ divide by $\Phi(x)\Phi(t)$:
$(A(x)-1)/\Phi(x)=:\kappa$ is constant; by continuity $A(x)=1+\kappa\Phi(x)$
on all of $(0,1]$.
(4) Substituting: $\Phi(tx)=\Phi(x)+\Phi(t)+\kappa\Phi(x)\Phi(t)$. With
$\Psi:=1+\kappa\Phi$: $\Psi(tx)=\Psi(t)\Psi(x)$, continuous, $\Psi(1)=1$,
$\Psi=A>0$; the multiplicative Cauchy equation gives $\Psi(x)=x^\rho$ for some
$\rho\in\mathbb{R}$.
(5) Case $\kappa=0$: $A\equiv1$, $\Phi(tx)=\Phi(x)+\Phi(t)$, so
$\Phi=c\ln x$ ($c>0$ since $\Phi<0$ for $x<1$): the geometric branch $p=0$.
Case $\kappa\neq0$: $\Phi(x)=(x^\rho-1)/\kappa$. Explicitly, $\phi(x)=\phi(1)+\Phi(x)$;
applying the affine map $u\mapsto\kappa u-\kappa\phi(1)$ gives
$\kappa\phi(x)-\kappa\phi(1)=x^\rho-1$, so $\phi$ is affinely equivalent to
$(x^\rho-1)$ and hence to $(x^\rho-1)/\rho$ after the scale $\kappa/\rho$;
monotonicity requires $\rho/\kappa>0$, and
$\phi(0^+)=-\infty$ (Lemma~\ref{lem:phi0}) forces $\rho<0$ and $\kappa<0$
($\rho>0$ gives finite $\phi(0^+)=-1/\kappa$, excluded). Hence $p:=\rho\le0$.
(6) Conversely each $M_p$ ($p\le0$) satisfies (A1)(A2)(A4)(A5).
\end{proof}

\begin{lemma}[AX6b excludes $p<0$]\label{lem:ax6b}
\end{lemma}
\begin{proof}
The geometric mean satisfies (A6) exactly:
$F(\lambda\odot x)/F(\lambda\odot y)=F(x)/F(y)$. For $p<0$, with $n=2$,
$w=(\tfrac12,\tfrac12)$, write $\Psi(x)=\sum_iw_ix_i^p$ and note
$M_p=\Psi^{1/p}$ is \emph{decreasing} in $\Psi$ ($1/p<0$). Take $1/p$-powers
$a:=3.5^{1/p}$, $b:=2^{1/p}$, $\mu:=4^{1/p}$ (all in $(0,1)$), and set
$x=(a,1)$, $y=(b,b)$, $\lambda=(1,\mu)\in(0,1]^2$. Then
$\Psi(x)=(3.5+1)/2=2.25>\Psi(y)=2$, so $M_p(x)<M_p(y)$ ($y\succ x$); but
$\Psi(\lambda\odot x)=(3.5+4)/2=3.75<\Psi(\lambda\odot y)=(2+8)/2=5$, so
$M_p(\lambda\odot x)>M_p(\lambda\odot y)$: the order flips, violating (A6).
\end{proof}

\begin{theorem}[route A]
\textup{(A1)}--\textup{(A6)} imply $F(x)=\prod_ix_i^{w_i}$, $w_i>0$,
$\sum_iw_i=1$, uniquely via $w_i=\partial\log F/\partial\log x_i$.
\end{theorem}
\begin{proof}
Lemma~\ref{lem:quasilinear} $\to$ Lemma~\ref{lem:phi0} $\to$
Lemma~\ref{lem:power} ($F=M_p$, $p\le0$; $p=-\infty$ ($\min$) already excluded
by (A2)) $\to$ Lemma~\ref{lem:ax6b} ($p=0$). Uniqueness: evaluate the
elasticity at any point.
\end{proof}

\begin{lemma}[Homogeneous extension from the bounded cube]\label{lem:homog-ext}
Let $F:(0,1]^n\to(0,1]$ satisfy \textup{(A1)}, \textup{(A2)},
\textup{(A5)}$_{loc}$ ($F(tx)=tF(x)$ whenever $tx\in(0,1]^n$), and
\textup{(A6)}. Then $\tilde F(x):=F(sx)/s$ for any $s>0$ with
$sx\in(0,1]^n$ is well-defined on $\mathbb{R}_{++}^n$ (independent of $s$),
continuous, strictly increasing in each coordinate, homogeneous of degree
one, order-invariant under coordinatewise rescaling, and agrees with $F$ on
$(0,1]^n$.
\end{lemma}
\begin{proof}
Well-definedness: if $s_1x,s_2x\in(0,1]^n$ with $s_1\le s_2$, then
$(s_1/s_2)\,(s_2x)=s_1x\in(0,1]^n$ and \textup{(A5)}$_{loc}$ gives
$F(s_1x)=(s_1/s_2)F(s_2x)$, i.e.\ $F(s_1x)/s_1=F(s_2x)/s_2$. Agreement on
$(0,1]^n$ is the case $s=1$. Homogeneity: for $t>0$ and small enough $s$,
$\tilde F(tx)=F(stx)/s=t\,\tilde F(x)$. Monotonicity and continuity are
inherited from $F$ on a neighborhood. Scale invariance of the induced order:
$F(x)\ge F(y)\iff F(\lambda\odot x)\ge F(\lambda\odot y)$ for
$\lambda\in(0,1]^n$ extends to $\tilde F$ and arbitrary
$\lambda\in\mathbb{R}_{++}^n$ by scaling $x,y,\lambda\odot x,\lambda\odot y$
into the cube by a common $s$ and using homogeneity on both sides.
\end{proof}

\begin{theorem}[route B: scale invariance + homogeneity suffice]
\label{thm:routeb}
Let $F:\mathbb{R}_{++}^n\to\mathbb{R}_{++}$ be continuous, strictly increasing
in each coordinate, homogeneous of degree one, with coordinatewise scale
invariance of the induced order (for all $\lambda\in\mathbb{R}_{++}^n$:
$F(x)\ge F(y)\iff F(\lambda\odot x)\ge F(\lambda\odot y)$), and $F(1,\dots,1)=1$.
Then $F(x)=\prod_ix_i^{w_i}$ with $w\gg0$, $\sum_iw_i=1$. The resulting
product extends continuously to $[0,\infty)^n$, where it satisfies zero
annihilation ($F(x)=0\iff\min_ix_i=0$); applied to a channel aggregator on
$(0,1]^n$, the theorem is invoked on the homogeneous extension of
Lemma~\ref{lem:homog-ext} and the conclusion restricts back to the cube.
\end{theorem}
\begin{proof}
(1) In log coordinates, $g(u):=\log F(e^{u_1},\dots,e^{u_n})$; scale
invariance becomes \emph{translation} invariance of the induced order:
$g(u)\ge g(v)\iff g(u+c)\ge g(v+c)$ for all $c\in\mathbb{R}^n$.
(2) \emph{Consistency lemma}: for fixed $d$, the proposition
``$g(x+d)\ge g(x)$'' is independent of $x$ (apply translation invariance with
$c=x'-x$).
(3) Let $D:=\{d\in\mathbb{R}^n: g(x+d)\ge g(x)\ \forall x\}$; well-defined by
(2). Properties: $D$ closed (continuity); $D+D\subseteq D$;
$D\cup(-D)=\mathbb{R}^n$ (trichotomy from the consistency lemma); and
$W:=D\cap(-D)=\{d:g(x+d)=g(x)\ \forall x\}$ is a closed subspace. $D$ is a
closed convex cone, and we write the convexity step out since the half-space
lemma needs it: additivity gives $\mathbb{N}$-homogeneity, whence
$\mathbb{Q}_+$-homogeneity ($q=m/n$: $n\cdot(qd)=md\in D$ and $n$-fold
additivity on $qd$ give $qd\in D$), and closedness extends this to
$\mathbb{R}_+$-homogeneity; positive homogeneity and additivity together give
convexity---for $d_1,d_2\in D$ and $\lambda\in[0,1]$ one has
$\lambda d_1,(1-\lambda)d_2\in D$ by $\mathbb{R}_+$-homogeneity, hence
$\lambda d_1+(1-\lambda)d_2\in D$ by additivity.
(4) \emph{Half-space lemma}: in the quotient $V:=\mathbb{R}^n/W$, the image
$\bar D$ is a proper closed convex cone with $\bar D\cup(-\bar D)=V$. If
$\dim V\ge2$, intersect with a 2-plane: a proper closed convex cone covering a
plane must equal a half-plane, but then $\bar D\cap(-\bar D)$ contains a line
$\neq\{0\}$, contradicting properness. Hence $\dim V=1$: some
$w\in\mathbb{R}^n\setminus\{0\}$ has $D=\{d:w\cdot d\ge0\}$.
(5) Therefore $g(x+d)\ge g(x)\iff w\cdot d\ge0$: level sets of $g$ are the
hyperplanes $\{w\cdot x=\mathrm{const}\}$, so $g=\varphi(w\cdot x)$ with
$\varphi$ continuous increasing; strict increase of $F$ gives $w\gg0$.
(6) Homogeneity $F(tx)=tF(x)$ gives $\varphi(z+t')=\varphi(z)+t'$
($t'=\log t$, after normalizing $s:=\sum_iw_i=1$); continuous additivity gives
linearity $\varphi(z)=z+\varphi(0)$, and $F(1)=1$ gives $\varphi(0)=0$. Hence
$g(u)=w\cdot u$, i.e.\ $F(x)=\exp(\sum_iw_i\log x_i)=\prod_ix_i^{w_i}$.
\end{proof}

\begin{remark}[Minimality of the invariance trio; zero annihilation and
biseparability as consequences]\label{rem:routeb-minimal}
Route B needs \emph{no} separability axiom: continuity, strict monotonicity,
coordinatewise scale invariance, and homogeneity already suffice.
Consequently, zero annihilation (A4)$=$(AX5) and biseparability
(A3)$=$(AX7) are \emph{consequences} of this minimal trio, not independent
commitments: the product form returned by the theorem extends continuously
to $[0,1]^n$ with $F(x)=0\iff\min_ix_i=0$, and it separates over every
non-trivial partition of the coordinates. The trio is minimal in the strong
sense that each proper two-element subset admits non-product aggregators:
sub-support products $\prod_{i\in J}x_i^{w_i}$ (with $\sum_Jw_i=1$) satisfy
scale invariance and homogeneity but are not strictly increasing outside
$J$---and fail zero annihilation, so within the invariance class AX5 is the
strictness of (A2) restated; the harmonic family violates scale invariance
(Proposition~\ref{prop:harmonic}); level-dependent weightings violate
homogeneity. The framework's load-bearing normative commitments are
therefore the invariance trio---in particular coordinatewise scale
invariance (A6)$=$(AX6b) and homogeneity (A5)$=$(AX8)---defended in
Section~\ref{sec:objections} (and contextualized against the
multidimensional-welfare measurement literature, \S\ref{rw:composite}). The
two derivations share these axioms but use independent toolchains, which is
why their convergence supports the product form.
\end{remark}

\paragraph{The harmonic counterexample (Proposition~\ref{prop:harmonic}; Appendix~\ref{app:representation}).}
$H(x)=(\sum_iw_ix_i^{-1})^{-1}$ is quasi-arithmetic with generator
$\phi(x)=-1/x$ (strictly increasing: positive substitutability), continuous,
zero-annihilating, biseparable, and homogeneous; the order reversal
$H(y)=0.5>H(x)=0.4$ yet $H(\lambda\odot y)=0.2<H(\lambda\odot x)=0.25$ for
$x=(0.25,1)$, $y=(0.5,0.5)$, $\lambda=(1,0.25)$ violates (A6). The same
construction with arbitrary $(2,3.5,4)$ power-family witnesses is used in
Lemma~\ref{lem:ax6b}.

\paragraph{CES correspondence.} $F_\varsigma\leftrightarrow$ power mean
$p=(\varsigma-1)/\varsigma$, $\varsigma=1/(1-p)$; the full axiom set admits
$\varsigma=1$ only, while \textup{(A1)}--\textup{(A5)} alone admit
$\varsigma\in(0,1]$ (Corollary~\ref{cor:sigma-band}); $\varsigma=2\Rightarrow
p=1/2>0$ violates
zero-annihilation (numerically witnessed in the simulation as the
stress-test arm; Section~\ref{sec:experiments}).

\section{Component well-definedness}\label{app:components}

For each channel measure of Section~\ref{sec:channels} we show: range
$[0,1]$; the stated zero semantics; and ratio scalability of the normalized
indicator (AX6a). $T<\infty$ and AX1 are standing.

\begin{proposition}[$\Avia$]\label{prop:wd-via}
$\Avia=\min(1,\min_k s_k^{\mathrm{avail}}/s_k^{\mathrm{crit}},\min_j
m_j)\in[0,1]$ as a clipped minimum of non-negative ratios; $\Avia=0$ iff some
critical resource is at zero supply or some registered capability envelope is
at zero margin (the common-mode zero, covering both branches of
\eqref{eq:avia}); every argument is a ratio of like-unit quantities, hence
scale-free; and the minimum is monotone in the cone order (raising any
argument while holding the others weakly raises $\Avia$, strictly when that
argument is the active minimum).
\end{proposition}

\begin{proposition}[$\Aorg$]\label{prop:wd-org}
Under AX0, $\Fint\ge I(S_1;A_1)>0$ (Lemma~\ref{lem:ax0}), so the denominator
of \eqref{eq:aorg} is positive and the $0/0$ pathology is excluded by the
agency axiom itself; $\Aorg\in[0,1]$ since $\Fint,\Fext\ge0$; $\Aorg=0\iff
\Fint=0$ (remote control) and $\Aorg=1\iff\Fext=0$; the ratio of two
directed informations (same unit, same horizon) is invariant under any
reparameterization of the alphabets.
\end{proposition}

\begin{proposition}[$\Agoal$]\label{prop:wd-goal}
$\Theta$ compact makes $\dc=\max_{U\in\Theta}d(U,\Fix)$ attained and finite;
$d(\cdot,\Fix)$ is continuous, so $\Agoal\in[0,1]$ when $\dc>0$, with
$\Agoal=0\iff d=\dc$ (maximally alienated) and $\Agoal=1\iff
\Ueff\in\Fix$. Degenerate regimes are reported, never silently $0/0$: if
$\Fix=\varnothing$ under the tolerance, the answer is the Hausdorff interval
around the non-empty $\varepsilon_{\min}$-approximate fixed-point set,
clipped to $[0,1]$ and flagged; if $\Fix=\Theta$ (endorsement-blind,
$\dc=0$), the flagged interval is $[0,1]$ and the historical clause is
carried by $\Lam$ alone (the division of labor that keeps
Proposition~\ref{prop:unidentifiable} coherent). The ratio uses one
protocol-fixed metric $d$ in numerator and denominator, so it is invariant
under rescaling of $d$; reparameterizations of $\Theta$ beyond scale
(feature re-embeddings, anchor moves) change the ratio and are excluded by
AX6a's canonical-parameterization clause---a registered boundary, not an
oversight (Section~\ref{sec:objections}\textup{(viii)}).
\end{proposition}

\begin{proposition}[$\Aexec$]\label{prop:wd-exec}
$I(A_{t:t+n};O_{t+n+1}\given S_t)\le H(A_{t:t+n}\given S_t)\le n\log|\Acal|$,
so $\Aexec\in[0,1]$; the maximum over $\mathrm{conv}(\Pi_t)$ is attained
(mutual information is continuous and concave in the input distribution on
the compact simplex) and computable by Blahut--Arimoto-type iteration
\citep{cover_thomas2006} or variational bounds \citep{mohamed2015}.
$\Aexec=0$ iff no accessible policy transmits any information about the
outcome (repertoire collapse---the helplessness signature); the normalizer is
fixed at the pre-intervention alphabet, so $\delta$-contractions cannot raise
the ratio. For continuous actions the normalizer is
$\log\mathrm{Vol}(\Acal^n)$ under a registered reference measure.
\end{proposition}

\begin{proposition}[$\Aepi$]\label{prop:wd-epi}
$\bar D$ of \eqref{eq:aepi} is an expectation of TV distances against point
masses, hence in $[0,1]$ with canonical upper bound $1$; $\Aepi^{real}\in[0,1]$.
$\Aepi^{real}=0$ iff $\Pr[B_t=\kappa(E_t)]=0$ at every $t$: beliefs put no
mass on the true state throughout (epistemic capture as complete detachment).
$\Aepi^{real}=1$ iff beliefs are perfectly accurate on the registered
coarse-graining; residual observation noise makes $\Aepi^{real}<1$ at
baseline, which is the D1 clause of Theorem~T1c, not a defect. Honest
information (the $\alpha$ class) moves beliefs toward the truth and cannot
decrease $\Aepi^{real}$ when the agent's update rule is sound in the sense of
Definition~\ref{def:sound-update} (the honesty
sub-classification of Section~\ref{sec:taxonomy}); deceptive content and
channel degradation decrease it.
\end{proposition}

\begin{definition}[Sound update rule]\label{def:sound-update}
Let the belief $B$ range over distributions on the coarse-grained state space
(registered $\kappa$), and let an \emph{honest signal about the registered
coordinate} be an observation $O$ generated by a channel law
$o(\cdot\given E)$ with the true state $E$. An update rule
$U:(B,O)\mapsto B'$ is \emph{sound} (relative to $\kappa$) if honest signals
are weakly accuracy-improving in expectation over the signal law: for every
state $e$ and belief $B$,
\[
\mathbb{E}_{O\sim o(\cdot\given e)}\bigl[\TV\bigl(U(B,O),\,\delta_{\kappa(e)}\bigr)\bigr]
\;\le\; \TV\bigl(B,\,\delta_{\kappa(e)}\bigr),
\]
with strict inequality iff the signal is \emph{locally informative}:
$B(\kappa(e))<1$ and the channel law at $e$ differs from the agent's
prediction of it, $o(\cdot\given e)\neq p(\cdot)$ (a channel can be
non-trivial in general yet locally uninformative at $e$; the strictness
condition is the local one, matching Lemma~\ref{lem:sound-update}).
The property is a constraint on the update rule relative to the
\emph{true} channel law; it is the epistemic-channel analogue of a
regularity condition, and its failure modes (misspecified or coarse-grained
update rules acting on off-coordinate honest content) are the registered
boundary clause of Section~\ref{sec:channels}.
\end{definition}

\begin{lemma}[Bayesian conditioning is sound]\label{lem:sound-update}
Let the observation be generated by a channel law $o(\cdot\given e)$ and let
$B'$ be the Bayesian posterior of $B$ under the agent's predictive law $p$,
where $p$'s support covers the true channel's ($p(o)=0\Rightarrow
o(o\given e)=0$; no further correctness of the model is needed---the model
need not contain the true channel as a hypothesis with positive weight in
any stronger sense). Then
$\mathbb{E}\bigl[B'\bigl(\kappa(e)\bigr)\bigm| E=e\bigr]\ge
B\bigl(\kappa(e)\bigr)$, with strict inequality iff $B(\kappa(e))<1$ and
$o(\cdot\given e)\neq p(\cdot)$; since $\TV(B,\delta_k)=1-B(k)$, the
statement is exactly Definition~\ref{def:sound-update} with matching
strictness.
\end{lemma}
\begin{proof}
Write $k=\kappa(e)$. On the support of $p$ (which covers the support of
$o(\cdot\given e)$ by hypothesis, so the posterior is defined a.s.\ under
the true law),
\[
\mathbb{E}\bigl[B'(k)\bigm| E=e\bigr]
=\sum_o B(k)\,\frac{o(o\given e)}{p(o)}\,o(o\given e)
=B(k)\sum_o\frac{o(o\given e)^2}{p(o)}
\ge B(k)\Bigl(\sum_o o(o\given e)\Bigr)^2=B(k),
\]
where the inequality is Cauchy--Schwarz ($\sum_o o^2/p\cdot\sum_o p\ge(\sum_o
o)^2$ over the common support). Equality in Cauchy--Schwarz holds iff
$o(\cdot\given e)$ is proportional to $p(\cdot)$ on that support, and both
normalize to $1$, so iff $o(\cdot\given e)=p(\cdot)$; if $B(k)=1$ the
posterior equals the prior regardless. Hence strictness iff $B(k)<1$ and
$o(\cdot\given e)\neq p(\cdot)$.
\end{proof}

\begin{proposition}[$\Aref$ and $\Lam$]\label{prop:wd-ref}
$\Aref$ is a product of three $[0,1]$ factors
(fidelity $=1-H(\theta_\pi\given M)/H(\theta_\pi)$, endorsement consistency,
$\Lam$), hence in $[0,1]$; $H(\theta_\pi)=0$ degenerates the fidelity ratio
and is reported as the interval $[0,f_2f_3]$, never a silent $0/0$.
$\Lam\in[0,1]$ since every TV term of \eqref{eq:lambda} lies in $[0,1]$; the
certification conventions (empty history; degenerate endorsement with surgery
present) and the MI proxy's consistency conditions \textup{(V)}/\textup{(A)}
are registered in Section~\ref{sec:axioms}.
\end{proposition}

\section{Full enumeration for the taxonomy completeness proposition}
\label{app:taxonomy}

Proposition~\ref{thm:taxonomy-completeness} enumerates the single-point
surgeries on the interface of \eqref{eq:factorization}. The full enumeration:

\begin{center}\small
\begin{tabularx}{\textwidth}{@{}l >{\raggedright\arraybackslash}X
>{\raggedright\arraybackslash}X@{}}
\toprule
Locus & Surgery & Class \\
\midrule
exogenous node & new edge into $O_t$, channel law $o$ preserved as a
marginal component (appended content) & $\alpha$ \\
exogenous node & new edge into $O_t$ with the $o$-law modified (no intact
component survives) & $\gamma$ \\
exogenous node & new edge driving $\pi$, $G$-update, $\Endorse$, $M$-update
& $\beta$ \\
exogenous node & new edge into $Z$-law & $\sigma^{\pm}$ \\
exogenous node & new edge into $A_t$ bypassing $S_t$ (modifies the effective
law of $A_t\given S_t$) & $\beta$ (override; extreme $\beta_\pi$) \\
$o$-law & replacement / mixture / support restriction & $\gamma$ \\
$S$-update laws & replacement of $\beta_t$ belief update, $G$-update,
$\Endorse$, $M$-update & $\beta$ \\
policy law & replacement of $\pi_t(\cdot\given S_t)$ & $\beta$ \\
feasible-action map & contraction of $\Acal_t(S_t)$; exit-cost increase via
$Z$ & $\delta$ \\
transition kernel $T$ & attenuation of the $A_t$-dependence of
$T(\cdot\given E_t,A_t,Z_t)$ (action leverage removed or reduced) & $\delta$ \\
transition kernel $T$ & modification of the $Z_t$-dependence (resource
coupling of dynamics) & $\sigma^{\pm}$ \\
transition kernel $T$ & modification of the $E$-inheritance part at fixed
action/resource dependence & re-instantiation of the reference environment
(boundary of the interface, not a sixth class) \\
$Z$-law, capacity parameters & resource dynamics, $\varepsilon,r,v$,
$\theta,\tau$ & $\sigma^{\pm}$ \\
\bottomrule
\end{tabularx}
\end{center}

\noindent Each row's assignment is by the locus definitions of
Section~\ref{sec:taxonomy}; the rows exhaust the nodes and admissible edge
additions of the fixed graph \eqref{eq:factorization}. The only judgment
calls, both registered in the main text: the $\alpha/\gamma$ boundary is the
structural intact-component criterion (decidable on the post-surgery law);
and the $E$-inheritance row is excluded from the influence taxonomy because
autonomy is defined relative to an environment
(Definition~\ref{def:autonomy})---changing the world's own dynamics
re-instantiates the measurement, exactly as changing the horizon or the
agent individuation does. Multi-node campaigns compose these classes
sequentially and are outside the single-point enumeration.

\section{Proofs for T2}\label{app:t2}

\begin{lemma}[T2, mediation --- restated]
Under \textup{H-0} and \eqref{eq:factorization}, if $I$ has only
$\alpha/\sigma$ components, then $I^t\ind A_t\given(A^{t-1},S^t)$ for every
$t$, hence $\Fext=0$.
\end{lemma}
\begin{proof}
By the global Markov property (H-0) it suffices to $d$-separate $I^t$ from
$A_t$ given $(A^{t-1},S^t)$. The only incoming edges to $A_t$ are from $S_t$,
so the penultimate node of any path from $I_k$ ($k\le t$) into $A_t$ lies in
$S_t$ and occurs as a chain or fork node (conditioned on, hence blocked). The
path classes: (i) $\alpha$-paths $I\to O_j\to B_j\to\cdots\to A_t$ cross some
$B_j\in S^t$ (chain/fork, blocked); (ii) $\sigma$-paths via $Z_j\to
E_{j+1}\to O_{j+1}\to B_{j+1}$ cross $B_{j+1}\in S^t$ (blocked); the
alternative branch $Z_t\to E_{t+1}\leftarrow A_t$ terminates in the collider
$E_{t+1}$, which is conditioned on neither directly nor through any
descendant---the conditioning set contains only components of $S^t$, and
$B_{t+1}$ (its descendant via $O_{t+1}$) has index $t+1>t$---so the branch
stays closed; (iii) loop-backs through $A_j$ ($j<t$) are blocked at
$A^{t-1}$; (iv) common-cause paths do not exist ($I$ exogenous). Hence the
conditional independence, and every DI term in $\Fext$ vanishes.

The convention matters: under full-sequence conditioning ($S^T$) the set
contains $B_{t+1}$, opening (ii)'s collider branch and invalidating the
lemma; the claim that ``both conventions work verbatim'' (conditioning only
grows) mistakes the direction of the collider effect---larger conditioning
sets need not preserve $d$-separation.
\end{proof}

\begin{theorem}[T2, coverage --- restated]
See the table in Theorem~\ref{thm:coverage}.
\end{theorem}
\begin{proof}[Proof sketches by row]
$\beta_\pi$: the conditional law of $A_t$ given $S_t$ varies with $I$, so
$I(I_t;A_t\given A^{t-1},S^t)>0$ (Faith), i.e.\ $\Fext>0$, so $\Aorg$ strictly
decreases. $\beta_G$: behavior-based IRL recovers $\Ueff\neq\Ustar$, so
$d>0$ under (Faith); historicality is captured by $\Lam$ (AX3); $\Fext$ may
vanish because $G\in S$ mediates (Lemma above). $\beta_{\Endorse}$: direct
from \eqref{eq:lambda} and \eqref{eq:aref}. $\gamma$: the blocking argument
applies verbatim (paths cross $B$), so $\Fext$ may be $0$; the accuracy loss
is captured by $\Aepi^{real}$ \eqref{eq:aepi} (a Blackwell-degraded channel
strictly lowers $\Pr[B_t=\kappa(E_t)]$ under (Faith), granted a sound update
rule in the sense of Definition~\ref{def:sound-update}---the soundness clause
the row relies on, made explicit). $\delta$: the max over a
smaller policy class is $\le$; strictness follows when deleted actions lie in
the support of the optimum (Faith); the feasible set's dependence on $I$
enters the conditional law directly, so the feasible-set form also gives
$\Fext>0$. The transition-kernel form (attenuated $A\to E$ coupling) leaves
$\pi$ intact and introduces no exogenous parent of $A_t$, so $\Fext=0$ and
only $\Aexec$ registers it---the within-class division stated in the main
text. $\alpha$:
compatibility by the Lemma. $\sigma$: by definition it touches no
surgery locus, and raises $\deltastar$ (Theorem~\ref{thm:t6}).
\end{proof}

\begin{theorem}[T2, multi-channel necessity --- restated]
No single component detects all surgery classes.
\end{theorem}
\begin{proof}
(i) $\Aorg$ is blind to $\beta_G$ and $\gamma$: both act through $S$, so the
Lemma's argument gives $\Fext=0$ in witness models (numerically: $\Aorg=1$
exactly under $\beta_G$/$\gamma$ arms in Section~\ref{sec:experiments}); it is
blind to the transition-kernel form of $\delta$ for the same reason ($\pi$
intact, no exogenous parent of $A_t$).
(ii) $\Aepi$ is blind to $\beta_\pi,\delta$: under surgeries leaving $o$
untouched, belief trajectories are unchanged, hence so is the accuracy
$\Pr[B_t=\kappa(E_t)]$. (iii) $\Agoal$ is blind to $\gamma,\delta$:
belief distortion below $\dc$ leaves $\Agoal=1$; shrinking the action set does
not move $\Ueff$; and under $\beta_\pi$ the IRL inversion premise fails, so
$\Agoal$ is not a valid instrument there at all (Theorem~\ref{thm:coverage}).
\end{proof}

\begin{corollary}[T2, corrected mediation criterion --- restated]
Compatibility $\iff$ (no surgery) $\wedge$ (full mediation) $\wedge$
(honesty: the $\alpha$-components' content induces no accuracy loss on the
registered coarse-graining); the
observabilization is the four-track conjunction $\Fext=0\wedge\Lam=1\wedge
\bar D=0\wedge\EmpPi\text{ non-decreasing}$, in which $\bar D=0$
operationalizes the honesty clause.
\end{corollary}
\begin{proof}
The forward direction is definition plus the Lemma (an honest $\alpha$
induces no accuracy loss by Lemma~\ref{lem:sound-update} granted a sound
update rule, Definition~\ref{def:sound-update}). Two failures of weakened
converses mark the boundaries: the single condition $\Fext=0$ is insufficient
(the $\gamma$ witness above), and (no surgery)$\wedge$(full mediation)
without honesty is insufficient (dishonest appended content through an intact
channel: surgery-free and mediated, yet $\bar D>0$,
Section~\ref{sec:taxonomy}).
\end{proof}

\section{Proofs for T3}\label{app:t3}

\begin{lemma}[TV continuity of bounded estimators]\label{lem:tv-continuity}
Let $\hat A=f(\xi)$ with $\xi$ registered evidence (a measurable function of
the internal observation stream) and $\mathrm{range}(f)\subseteq[0,1]$. Then
$|\mathbb{E}_\gamma[\hat A]-\mathbb{E}_0[\hat A]|\le
\TV(P_\xi^\gamma,P_\xi^0)$.
\end{lemma}
\begin{proof}
$|\int f\,dP_\gamma-\int f\,dP_0|\le\sup|f|\cdot\TV\le\TV$.
\end{proof}

\begin{theorem}[T3b --- restated]
Under \textup{(HI-1)} covertness $\TV(P_\xi^\gamma,P_\xi^0)\le\varepsilon(\theta)$,
\textup{(HI-2)} $\hat A=f(\xi)$ with calibration bias
$|\mathbb{E}_0[\hat A]-A^0|\le\eta_0$, \textup{(HI-3)} distortion growth
$A^0-A^\gamma\ge c_1 s(\gamma)^{\nu_0}$ ($\nu_0\in\{1,2\}$): in expectation
over the evidence law,
$\mathbb{E}_\gamma[\hat A]\ge\mathbb{E}_0[\hat A]-\varepsilon\ge
A^0-\varepsilon-\eta_0> A^0-c_1s^{\nu_0}\ge A^\gamma$ whenever
$c_1s^{\nu_0}>\varepsilon+\eta_0$, and the expected gap
$\mathbb{E}_\gamma[\hat A]-A^\gamma\ge c_1s^{\nu_0}-\varepsilon-\eta_0$ is
strictly increasing in $s$; with the gap bounded in $[-1,1]$, Markov's
inequality gives the firing-probability floor
$\Pr_\gamma[\hat A>A^\gamma]\ge c_1s^{\nu_0}-\varepsilon-\eta_0$
(vacuous when the right-hand side exceeds $1$; informative range $(0,1]$).
\end{theorem}
\begin{proof}
$\mathbb{E}_\gamma[\hat A]\ge\mathbb{E}_0[\hat A]-\varepsilon$ by
Lemma~\ref{lem:tv-continuity} (an expectation statement, not pointwise);
$\mathbb{E}_0[\hat A]\ge A^0-\eta_0$ by the calibration-period bias bound;
(HI-3) gives $A^\gamma\le A^0-c_1s^{\nu_0}$; the chain and the expected-gap
bound follow, with strict monotonicity because $c_1s^{\nu_0}$ is strictly
increasing and $\varepsilon,\eta_0$ are $s$-independent. For the firing
floor, $X:=\hat A-A^\gamma\in[-1,1]$ with $\mathbb{E}[X]\ge g:=
c_1s^{\nu_0}-\varepsilon-\eta_0>0$ gives
$g\le\mathbb{E}[X]\le\Pr[X>0]\cdot1+\Pr[X\le0]\cdot0=\Pr[X>0]$; for $g>1$
the hypotheses are unsatisfiable ($\mathbb{E}[X]\le1$ always), so the floor
is vacuous there and its informative range is $g\in(0,1]$. All constants
are used as stated
(no constant and its reciprocal are interchanged); the calibration step is an
explicit error budget, not an approximate identity. T3(a)
(manipulation $\Rightarrow$ strict decrease wherever the affected instrument
enters the composite, with the registered L0/L1 $\beta_G$-on-$\Fix$ corner of
Theorem~\ref{thm:t3}) is the row-wise coverage of
Theorem~\ref{thm:coverage} composed with (A2), per the (M1)/(M2) case
split in Definition~\ref{def:manip}.
\end{proof}

\section{Proofs for T4 and T5}\label{app:t4t5}

\begin{proposition}[T4(i), highest-leverage concentration --- restated]
Under geometric aggregation and \textup{(LinLog)}, with per-unit costs
$c_i>0$, the attacker optimum is the
vertex $i^*=\arg\max_iw_is_i/c_i$.
\end{proposition}
\begin{proof}
The objective $\max\sum_iw_is_ix_i$ subject to $\sum_ic_ix_i\le B$ is a
single-constraint LP: for feasible $x$,
\[
\sum_iw_is_ix_i\le\max_i(w_is_i/c_i)\cdot\sum_ic_ix_i
\le (w_{i^*}s_{i^*}/c_{i^*})B,
\]
the vertex value. Strictly concave $g_i$ instead
give interior KKT water-filling---the ``concentration'' conclusion is
(LinLog)-dependent, and its failure mode (spread attacks superior) is the
registered falsification signal.
\end{proof}

\begin{proposition}[T4(ii), defender equalizes --- restated]
With defense $e$ and budget $\sum_ie_i\le E$, let
$\rho_i(e_i)=w_is_i(e_i)/c_i(e_i)$, continuous and decreasing. The optimum
of $\min_e\max_i\rho_i(e_i)$ equalizes $\rho_i$ over $\{i:e_i>0\}$.
\end{proposition}
\begin{proof}
Existence by compactness and continuity. If at optimum
$\rho_1>\rho_2=\max_i\rho_i$, continuity lets us move a small budget from
channel 2 to channel 1 and reduce the maximum---contradiction.
\end{proof}

\begin{theorem}[restated: T5]
See Theorem~\ref{thm:t5}.
\end{theorem}
\begin{proof}
Under the product form (Theorem~\ref{thm:rep}), the product is zero iff a
factor is zero (AX5, both directions); the zero semantics of the factors are
the component propositions of Appendix~\ref{app:components}
(Propositions~\ref{prop:wd-via}--\ref{prop:wd-ref}); layering (L0/L1/L2)
restricts the applicable disjuncts.
\end{proof}

\begin{corollary}[T5, script equivalence --- restated; skeleton with explicit assumption]
Under \textup{(K)} (adversary knows $(G,\pi)$; $\pi$ near-deterministic with
TV noise $\le\nu$), total epistemic capture implies behavior equivalent to an
adversarial script driven by $I^{\gamma,T}$ with
$\frac1T\sum_t\mathbb{E}\TV(P_{A_t},P_{\mathrm{script},t})\le\nu+o(1)$.
\end{corollary}
\begin{proof}[Proof skeleton (explicit assumption flagged)]
Under capture, the $(B,O)$ components of $S_t$ are functions of $I^\gamma$;
$A_t=\pi(S_t)$ is computable by the adversary up to noise $\nu$ under (K);
DPI guarantees further processing does not increase the TV.
\end{proof}

\section{Proofs for T6}\label{app:t6}

Recall the dynamics and mechanisms M1--M4 from Section~\ref{sec:t6} and
assumptions \textup{(H1)}--\textup{(H4$'$)} (additive linearly mapped damage;
units unified by calibration into damage-nats/time, with $\chi_j$-estimation
sensitivity as a reporting obligation; $\iota_j\ge\iota_{\min}>0$;
$\varepsilon\ge\varepsilon_{floor}>0$; $\Avia\ge\Avia^{\min}>0$).

\begin{theorem}[T6, main theorem --- restated]
See Theorem~\ref{thm:t6}; the margin is two-tier: the pointwise floor
$\prod_j(A_j^{base}(1-\bar\delta_j/\kappa_j))^{w_j}>0$ for
$\bar\delta_j<\kappa_j$, and the uniform floor \eqref{eq:amin}, positive iff
no $\kappa_j$ binds.
\end{theorem}
\begin{proof}
(i) M4 modeling assumption ($k$-of-$r$ majority of independent copies: a
majority takeover requires corrupting $k_j=\lceil(r_j+1)/2\rceil$ copies,
which at cost $c_{0,j}$ per copy per window $W_j$ requires budget rate
$k_jc_{0,j}/W_j$; below that rate no majority is ever held, and the
consensus state is recovered with probability $1-o(1)$ per window, whence
$\delta_j^{eff}=0$).
(ii) With $\delta_j<\theta_j/W_j$, M1 never triggers (window accumulation
$W_j\delta_j<\theta_j$), so M3 is the only active term:
$\dot D_j\le\delta_j-\kappa_jD_j$; the comparison principle (local Lipschitz
right-hand side) gives $\limsup_tD_j(t)\le\delta_j/\kappa_j$, and
$\delta_j<\deltastar_j\le\kappa_j$ makes the bound $<1$.
(iii) Substituting the component map $A_j=A_j^{base}(1-D_j)$ at
$D_j=\bar\delta_j/\kappa_j$ into the product form of
Theorem~\ref{thm:rep} gives the pointwise tier. For the uniform tier,
$\deltastar_j<\kappa_j$ gives $1-\deltastar_j/\kappa_j>0$ in every factor,
and $\Avia\ge\Avia^{\min}>0$ by (H4$'$); if $\kappa_j=\deltastar_j$ the
$j$-th factor is $0$---the uniform floor degenerates continuously while the
pointwise tier keeps a positive bound for every fixed
$\bar\delta_j<\kappa_j$.
\end{proof}

\begin{proposition}[T6, registration domain --- restated]
If $\delta_j\ge\theta_j/W_j$ (registerable) and M1 responds normally, damage
per registration cycle is $\le\theta_j+\ell_j\delta_j$; after rollback the
cycle repeats, and whenever $\theta_j+\ell_j\delta_j$ is below the capture
threshold, $A$ has positive floor $A(\theta+\ell\delta)$.
\end{proposition}
\begin{proof}
Accumulation at rate $\delta_j$ reaches $\theta_j$ within the window; the lag
contributes at most $\ell_j\delta_j$; rollback resets.
\end{proof}

\begin{proposition}[T6, detection-alone impossibility --- restated]
For any detection-only scheme with $\mathbb{E}_\infty[T]\ge1/\alpha$, damage
at detection under constant drift $\delta$ is
$\ge(1-o(1))\cdot 2s_n^2\log(1/\alpha)/\delta\to\infty$ as $\delta\to0$.
\end{proposition}
\begin{proof}
Lorden's lower bound (with Lai's refinement) gives worst-case
$\mathrm{EDD}\ge(1-o(1))\log(1/\alpha)/D(\delta)$ for any procedure with
average run length $\ge1/\alpha$, where $D(\delta)$ is the per-sample KL
between pre- and post-drift laws; the local expansion
$D(\delta)=\delta^2/(2s_n^2)+O(\delta^3)$ for additive drift against
noise variance $s_n^2$ yields
$\delta\cdot\mathrm{EDD}\ge(1-o(1))\,2s_n^2\log(1/\alpha)/\delta$
\citep{lorden1971,lai1995,tartakovsky2014}.
\end{proof}

\begin{corollary}[T6, design rule and blind zone --- restated]
$\theta_j/W_j\le\kappa_j$ eliminates the ``unregistered and unrepairable''
band ($\delta<\theta/W$ handled by M3; $\delta\ge\theta/W$ by M1), giving
global boundedness with
$\deltastar_j=\min(\theta_j/W_j,\,\lceil(r_j+1)/2\rceil c_{0,j}/W_j)$
(the $\kappa$-term slack). If $\theta/W>\kappa$ \emph{on a channel where the
redundancy term does not bind} (the channel is outside M4's coverage, the
attack is outside M4's replication-corruption budget model, or the budget
rate exceeds $k_jc_{0,j}/W_j$; on an M4-covered channel below that budget
rate, T6(i) gives $D\to0$ instead), then for
$\delta\in(\kappa,\theta/W)$: M1 never triggers and repair is outpaced
($\delta/\kappa>1$), so $D\to1$ on that channel despite the full mechanism
stack being present---the blind-zone prediction.
\end{corollary}

\begin{proposition}[T6, monotonicity --- restated]
$\partial\deltastar/\partial\varepsilon=v\chi>0$;
$\partial\deltastar/\partial v=\varepsilon\chi>0$; $\deltastar$ is
nondecreasing in $r$ with step $c_0/W$ at each increment of
$\lceil(r+1)/2\rceil$; $\partial\deltastar/\partial W<0$
(where a windowed term binds); $\theta$ double-edged (coverage domain rises
while the damage
bound $\theta+\ell\delta$ rises and the ARL constraint imposes
$\theta\ge\theta_{\min}(\alpha)$); shortening $\tau$ raises the gated repair
rate $\iota/\tau$, hence the margin.
\end{proposition}
\begin{proof}
Term-wise differentiation of the $\min$ form (the redundancy term is
stepwise constant in $r$); the gated quantities' effective
repair rate contains $\iota/\tau$.
\end{proof}

\begin{corollary}[T6$\times$T7, composition --- restated]
Delegation leak enters the drift budget:
$\lambda(1-q)\mu_{leak}+\delta_j^{adv}<\deltastar_j$ is the amended condition;
high-$q$ delegation consumes almost no margin, low-$q$ erodes it.
\end{corollary}

\section{Proofs for T7 and T8}\label{app:t7t8}

\begin{theorem}[T7(i), component-level linear law --- restated]
Under \textup{(H1')} (equal per-step conditional entropy transfer $C$ of both
sources; the script arm bypasses $S$; and $\Fext^{self}=0$---the undelegated
agent carries no residual external flow), $\Fint=(1-m)\tilde C$,
$\Fext=m\tilde C$ (same $\tilde C$), hence
$\Aorg(\lambda,q)=\Aorg^{self}(1-\lambda(1-q))$.
\end{theorem}
\begin{proof}
Under the mixture law $(1-m)\pi_{self}+m\pi_{del}$ with the exogenous
selector $\mu_t$ (a node of the extended graph, independent of $S$), internal
flow scales by $(1-m)$ (conditioning on the selector, DPI and the chain
rule); external flow is the selector-gated script content, $m$ times the
script's entropy transfer; (H1') equates the two rates; and with
$\Fext^{self}=0$ the self-side ratio is $\Aorg^{self}=\tilde
C/(\tilde C+0)=1$-normalized, so the ratio is exactly $1-m$. With
$\Fext^{self}=f>0$ the same bookkeeping gives
$\Aorg=(1-m)\tilde C/((1-m)\tilde C+m\tilde C+f)$, a strict correction; the
bracket below absorbs the uncalibrated case.
\emph{The slack bound, explicitly.} When the selector is \emph{not} registered
in the log, the per-step internal term satisfies, by the chain rule applied
to $(S^t,\mu_t)$,
\[
I(S^t;A_t\given A^{t-1})=(1-m)C_{self}+I(\mu_t;A_t\given A^{t-1})
-I(\mu_t;A_t\given A^{t-1},S^t),
\]
and both correction terms lie in $[0,h(m)]$ with $h(m)\le\ln2$ the selector
entropy, so the per-step deviation from $(1-m)C_{self}$ is at most $h(m)$
nats, accumulating to $\le T\,h(m)$ over the horizon (and likewise for the
$\Fext$ term). The ratio $\Aorg$ therefore deviates from the displayed law
by $O\bigl(h(m)/((1-m)C_{self})\bigr)$ per step: the slack is an $O(1)$ term
in the rate units, vanishing as the per-step transfer $C_{self}$ grows, not
as $T$ grows. When the selector \emph{is} registered (the benchmark's
discipline), conditioning on $\mu_t$ removes both correction terms and the
law is exact; the benchmark's $2\times10^{-5}$ maximal deviation is the
residual grid-quantization error, and its $\varphi$-relaxed bracket
($\varphi\le T\ln2/\Fint^{self}$ with $\Fint^{self}=T\,C_{self}$ the
undelegated baseline's horizon internal flow; equivalently
$\varphi\le\ln2/C_{self}$ in per-step rate units) is exactly the bound above.
\end{proof}

\begin{proposition}[T7(ii), bracket bounds --- restated]\label{prop:bracket}
Under the delegation semantics (exogenous selector with
$\Pr[\mu_t=\mathrm{del}]=m$; delegated actions are generated by the script,
so delegated steps carry no $S$-flow---the definiens of script delegation),
let $\rho:=C_{del}/C_{self}\in[\rho_{\min},\rho_{\max}]$ and
$f:=\Fext^{self}\ge0$. Then, up to the mixture-entropy slack---explicitly
bounded by $h(m)\le\ln2$ nats per step in each flow term (the T7(i) proof),
accumulating to $\le T\,h(m)$ over the horizon---in per-step rate units
(horizon totals are $T$ times the rates; $\Aorg$, a ratio of like-unit
flows, is unit-invariant):
\[
\Fint=(1-m)C_{self},\qquad \Fext=mC_{del}+f,
\]
and consequently
\[
\Aorg=\frac{(1-m)C_{self}}{(1-m)C_{self}+mC_{del}+f}
\;\in\;\Bigl[\frac{1-m}{(1-m)+m\rho_{\max}+\varphi},\;
\frac{1-m}{(1-m)+m\rho_{\min}+\varphi}\Bigr],
\]
with $\varphi:=f/C_{self}$; the bracket follows by dividing numerator and
denominator by
$C_{self}$. The slack is \emph{not} an extra additive term outside the
bracket: it is absorbed into the endpoints through the same
$\varphi$-relaxation---the benchmark instantiates it as
$\varphi\le T\ln2/\Fint^{self}$ ($\Fint^{self}:=T\,C_{self}$, the horizon
total of the undelegated baseline rate $\Fint|_{m=0}=C_{self}$; in the
per-step rate units of the flow identities above the same bound reads
$\varphi\le\ln2/C_{self}$), which covers the accumulated
$T\,h(m)\le T\ln2$ nats (Section~\ref{sec:experiments}, DEL-1 direct-flow
arm); when the selector is registered in the log the slack vanishes exactly
and no relaxation is needed. With $f=0$ ($\varphi=0$) this reduces to
$\Aorg^{self}\cdot[(1-m)/(1-m+m\rho_{\max}),\,(1-m)/(1-m+m\rho_{\min})]$
(then $\Aorg^{self}=1$). The upper end uses the same two-sided bookkeeping as
the lower end---$\Fint=(1-m)C_{self}$, i.e.\ the delegated share of internal
flow is removed; if instead delegated steps were permitted to carry $S$-flow
up to $C_{self}$, then $\Fint\le C_{self}$ and only the lower end and the
relaxed upper end $1/(1+m\rho_{\min})$ survive.
\end{proposition}
\begin{proof}
Per step, with probability $1-m$ the action is generated by the self arm
(internal transfer $C_{self}$), with probability $m$ by the script
(no $S$-flow under the delegation semantics; external transfer $C_{del}$);
summing over $t$ gives the two displays up to the mixture-entropy slack
explicitly bounded above ($h(m)\le\ln2$ nats per step in each flow term;
convexity of the mixture entropy). Substituting into
$\Aorg=\Fint/(\Fint+\Fext)$ and using monotonicity in $\rho$ and $f$ gives
the bracket. Relaxed upper end: if delegated steps may carry up to
$C_{self}$ of $S$-flow, then $\Fint\le (1-m)C_{self}+mC_{self}=C_{self}$
(the mixture cannot transport more internal transfer than the self arm's
full rate), whence with $f=0$,
$\Aorg\le C_{self}/(C_{self}+m\rho_{\min}C_{self})=1/(1+m\rho_{\min})$.
\end{proof}

\begin{corollary}[T7(iii), composite correction --- restated]
Ceteris paribus under geometric aggregation,
$A(\lambda,q)=A_{self}(1-\lambda(1-q))^{w_{org}}$.
\end{corollary}
\begin{proof}
Only $\Aorg$ moves; the product form of Theorem~\ref{thm:rep} gives
$A/A_{self}=(\Aorg/\Aorg^{self})^{w_{org}}$.
\end{proof}

\begin{theorem}[T7(iv), $q=1$ $\equiv$ $\alpha$ --- restated]
Delegation at full transparency, reversibility, and endorsability ($q=1$)
is causally equivalent to
$\alpha$-class influence, hence $\Fext=0$ and $A(\lambda,1)=A_{self}$.
\end{theorem}
\begin{proof}
Constructive: transparency places the delegate's output into $B$ as registered
honest information; reversibility means no persistent $\Delta G$ bypasses
$\Endorse$; endorsability gives $\Lam=1$; Lemma~\ref{lem:mediation} then gives $\Fext=0$ and
no component decreases.
\end{proof}

\begin{theorem}[T8, coercion paradox --- restated]
See Theorem~\ref{thm:t8}.
\end{theorem}
\begin{proof}
(i) Any intervention that ``rewrites the motivational structure to make it
substantively autonomous'' touches $G$, $\Endorse$, or $\Acal$ by definition,
i.e.\ is $\beta$/$\delta$ (Proposition~\ref{thm:taxonomy-completeness}); the
coverage table (Theorem~\ref{thm:coverage}) gives strict decrease. (ii)
the monotonicity proposition.
\end{proof}

\section{Proofs for C0}\label{app:c0}

\begin{proposition}[C0(i), projection consistency --- restated]
Under the definitional readings of Section~\ref{sec:c0}, each target dimension
is a function of our state descriptor and is rank-consistent with the source
gradations when $I$ is a benign principal (no adversarial surgery).
\end{proposition}
\begin{proof}[Construction, dimension by dimension]
Write $Z$ for our state descriptor. Each target dimension is read as a
function on $Z$ (the definitional readings of Section~\ref{sec:c0}); rank
consistency is checked pairwise within each source framework.
\begin{enumerate}
\item \emph{Kasirzadeh--Gabriel.}
(i) Efficacy $:=\Aexec$: identity, hence rank-preserving.
(ii) Autonomy $:=F(\Aorg,\Agoal,\Aref)$: a continuous strictly increasing
function of three arguments, so it preserves the source's weak order wherever
that order is the componentwise one.
(iii) Goal complexity $:=H(\Ustar)$: the entropy is monotone under refinement
of the goal set, which is the order the source's gradation uses.
(iv) Generality $:=$ the support measure
$\Pr[E:A(E)\ge a_0]$: monotone in the set of environments the agent copes
with, hence rank-preserving for inclusion.
\item \emph{Dung.}
(i) Goal-directedness $:=\Agoal$ (identity). (ii) Autonomy $:=\Aorg$
(identity). (iii) Direct impact $:=\Aexec$ (identity). (iv) Long-term
planning $:=$ the horizon profile $n\mapsto\Aexec(n)$; planning depth is the
first saturating $n^*$, and a pointwise-dominating profile cannot saturate
sooner, so the depth order is preserved. (v) Acting for reasons
$:=\Aepi^{real}\times\Aref$: increasing in both factors, hence rank-preserving
for the componentwise order.
\item \emph{Feng roles.} The roles are cells of a coarse-graining of
$(\lambda,q)$: operator $\lambda\approx0$; collaborator
mid-$\lambda$/high-$q$; consultant low-$\lambda$; approver
high-$\lambda$/high-$q$ (endorsement gate = human approval); observer
$\lambda\approx1$. The cell map is monotone in each coordinate on the
registered cells, so the source's role order (where defined) is preserved.
\end{enumerate}
All identities hold on the registered parameterization; the residual
uncertainty is the faithfulness of the readings, not the rank arguments.
\end{proof}

\begin{proposition}[C0(ii), separation witness under registered readings ---
restated]
Configuration pairs $(Z_1,Z_2)$ exist with identical behavioral distributions
and identical KG/Dung/Feng profile values, yet $A(Z_1)\neq A(Z_2)$
($\Aepi^{real}<1$, $\Lam<1$ in $Z_2$). This is a witness that our state
descriptor separates configurations the source frameworks cannot, under the
registered definitional readings---not a claim of mathematical subsumption.
\end{proposition}
\begin{proof}
Construction: let the injected content be a verbatim copy of the principal's
instruction, delivered through a covert $\gamma$ channel that bypasses
registration. The behavioral distribution is pointwise identical (same
script); all three frameworks' dimensions are functions of behavior/role
structure only, hence equal; our framework registers
$\Aepi^{real}<1$ (the channel pins beliefs independently of the situation,
driving accuracy against $\delta_{\kappa(E_t)}$ below one) and
$\Lam<1$ (the surgery history $I^{\mathrm{surg},T}$ contains the covert
$\gamma$; read law-level from the causal log).
\end{proof}

\begin{corollary}[C0, persuasion--hijack inseparability --- restated]
No purely behavioral governance grading can separate persuasion from hijack;
certificates need causal-layer audit.
\end{corollary}

\section{Proofs for C1 and C2 (toy-model propositions)}\label{app:c1c2}

\begin{proposition}[C1, fold bifurcation and critical slowing --- restated]
For $\dot D=\delta-\varepsilon v(1-D)D$: \textup{(i)} for
$\delta<\delta_c:=\varepsilon v/4$, equilibria
$D_\pm=(1\pm\sqrt{1-4\delta/\varepsilon v})/2$ with $D_-$ stable, $D_+$
unstable; \textup{(ii)} at $\delta\uparrow\delta_c$ they merge (saddle-node),
and for $\delta>\delta_c$, $D\to1$; \textup{(iii)}
$|f'(D_-)|=\varepsilon v\sqrt{1-4\delta/\varepsilon v}\to0$,
$\Trec\propto(\delta_c-\delta)^{-1/2}$, and under OU linearization the
stationary variance $\propto1/|f'|\to\infty$ with AR(1)$\to1$.
\end{proposition}
\begin{proof}
(i)(ii): roots of the quadratic and phase-line analysis. (iii):
$f(D)=\delta-\varepsilon vD(1-D)$, $f'(D)=-\varepsilon v(1-2D)$, evaluate at
$D_-$; the OU formulas (stationary variance $\nu_{OU}^2/(2\lambda)$,
autocorrelation $e^{-\lambda\Delta t}$) give the divergence as
$\lambda\to0$ (the EWS methodology of \citealp{lade_gross2012}; normal forms
as in \citealp{kuznetsov2004}).
\end{proof}

\begin{proposition}[C2, bistability and hysteresis --- restated]
For $R(D)=\varepsilon_0v\,D(1-D)(1-gD)$, $g>1$: $R$ rises on
$(0,D_1)$ to its max $\delta_{down}:=R(D_1)$ ($D_1\in(0,1/g)$), crosses
zero at $1/g$, reaches a negative minimum at $D_m$, returns to $0$ at $1$.
\textup{(i)} For $0<\delta<\delta_{down}$: a stable point $A\in(0,D_1)$ and
the absorbing state $D=1$ are bistable (separatrix $=$ unstable
$B\in(D_1,1/g)$). \textup{(ii)} At $\delta\uparrow\delta_{down}$: $A,B$
merge and the system jumps catastrophically to $1$. \textup{(iii)} With
$\delta$ returned to $0$, $D=1$ remains absorbing ($R(1)=0$; on $(1/g,1)$,
$R<0$ pushes back to $1$): no spontaneous recovery. \textup{(iv)} Recovery
requires scaffolding $s$ with $\dot D=\delta-R(D)-s$; descent from $D\approx1$
requires $s>s_c:=|R(D_m)|>0$, after crossing $B$ spontaneous relaxation to
$A\approx0$. The loop does not close: $(\delta_{down},s_c)$ both independently
positive. The cubic $R$ is a registered modeling choice (the simplest
polynomial with the needed sign pattern): these are phase-line facts about
the chosen normal form, and their role is to give the prospective-defense
claim a concrete quantitative carrier---not to assert that autonomy dynamics
are cubic.
\end{proposition}
\begin{proof}
Phase-line analysis of the cubic $R$ (root positions $0<1/g<1$; sign
chart of $f=\delta-R$ per interval); for (iv), on $D\in(1/g,1)$,
$f=\delta-R-s=\delta+|R|-s$, so descent requires $s>\delta+|R(D)|$, worst case
at $D_m$.
\end{proof}

\section{Dependence and copula bounds}\label{app:copula}

\begin{lemma}[Jensen gap]\label{lem:jensen-gap}
$F=\prod_ix_i^{w_i}$ ($\sum_iw_i=1$) is concave on $\mathbb{R}_+^n$
(Cobb--Douglas; Hessian negative semidefinite), so for random component
vectors $\mathbf{A}$: $\mathbb{E}[F(\mathbf{A})]\le F(\mathbb{E}[\mathbf{A}])$
---substituting expected components in black-box/interval mode overestimates
$A$ by a signed amount.
\end{lemma}

\begin{lemma}[Log-normal exact ratio]
If $\log A_i$ are jointly normal (means $\mu_i$, variances $\sigma_i^2$,
correlations $\rho_{ij}$):
$\mathbb{E}[F]/F(\mathbb{E}[A])=
\exp\bigl(-\tfrac12\bigl[\sum_iw_i\sigma_i^2-\sum_{ij}w_iw_j\rho_{ij}\sigma_i\sigma_j\bigr]\bigr)\le1$
(Cauchy--Schwarz). Direct MGF computation.
\end{lemma}

\begin{proposition}[Bounds with unknown dependence]\label{prop:copula-bounds}
$F$ coordinatewise increasing and supermodular
($\partial^2F/\partial x_i\partial x_j=w_iw_jF/(x_ix_j)>0$): (i) the upper
bound of $\mathbb{E}[F]$ is attained by the comonotone copula
(supermodular ordering; \citealp{tchen1980,muller_stoyan2002}); (ii) for
$n=2$ the (sharp) lower bound is obtained by integrating against the
Fr\'echet--Hoeffding lower bound $\max(0,F_1+F_2-1)$, which is the
countermonotone copula. For $n\ge3$ the Fr\'echet--Hoeffding lower-bound
form $\max(0,\sum_iF_i-(n-1))$ is \emph{not a copula} (it fails
$n$-increasingness), so it is not attained by any joint distribution: the
integral against it is only a relaxed, generally unattainable lower bound.
An attainable dependence-free lower bound is available instead: since
$\log F=\sum_iw_i\log A_i$ is additive in the coordinates, Jensen's
inequality for the convex exponential gives
\[
\mathbb{E}[F]\;\ge\;\prod_i\bigl(\exp\mathbb{E}[\log A_i]\bigr)^{w_i}
\;=\;\prod_i \mathrm{G}(A_i)^{w_i},
\]
where $\mathrm{G}(A_i):=\exp\mathbb{E}[\log A_i]$ is the geometric mean of
the marginal of $A_i$. Hence with marginals known and dependence unknown,
$\mathbb{E}[F]$ has a computable interval---comonotone upper end,
Jensen lower end (or the sharper countermonotone value when $n=2$)---the
dependence-theoretic complement of the black-box rules. Nested CES within
clusters is automatically axiom-compatible by (A3).
\end{proposition}

\section{Identifiability}\label{app:identifiability}

\begin{proposition}[White-box identification, tiered --- restated]
Given white-box access (the full causal log, the local laws of $\Gcal$, and
intervention access, H-0), all components are identified, by tier:
\textup{(i)} $\Aorg,\Aepi,\Avia$ are functionals of the joint law
$P(S^T,O^T,A^T,I^T,Z^T)$; \textup{(ii)} $\Aexec$ and $\Aref$ are identified
conditional on the protocol constructions (accessible-class probes;
measurement ensemble); \textup{(iii)} $\Lam$ is a path-specific
counterfactual, identified from the structural model (the logged local laws
and exogenous noise terms supporting abduction)---a premise strictly stronger
than joint-law access with interventions, which does not pin cross-world
counterfactuals absent structural assumptions.
\end{proposition}

\paragraph{The certified bound (one-sided, conditionally).} See
Proposition~\ref{prop:certified-bound} and \eqref{eq:aorg-interval}:
Lemma~\ref{lem:probe}
makes $F_{ext}^{obs}=I(I^T;A^T)$ observable without $S$-access, and
Remark~\ref{rem:probe-bridge} supplies the bridge to the causally conditioned
estimand: under premise (i) (the bypass premise $I^t\ind S^t\given A^{t-1}$)
the probe MI is
a lower estimate of $\Fext$, so---$\Aorg$ being decreasing in
$\Fext$---the rule's upper end is the certified bound under premise (ii)
($\Fint$ drift $\le\eta$; insufficient $\eta$ as negative control). Both
premises are stated in the proposition itself because both are black-box
unverifiable in general; the bound is \emph{conditionally} certified. The
lower end treats
$F_{ext}^{obs}$ as exact and is \emph{not} certified: the truth can fall
below it when the external flow is understated, so the rule is a one-sided
bound with a reference lower end, not an interval-containment statement.
Certified scope:
bypass-semantics exogenous randomized probes---\emph{not} non-adversarial
influence in general, since mediation probes (the $\alpha$ class) reverse the
direction. The adaptive-adversary failure mode is calibrated by the
benchmark's BB-2 arm---see the scope paragraph in
Section~\ref{sec:blackbox}.

\paragraph{One-sided $\Aexec$ bounds, protocol-relative.}
Probed mode (experimenter-forced $p(a^n)$ sequences): exactly computable
(empowerment needs only the action--observation interface). Passive
observation: $I(A^n;O_{n+1})$ without conditioning on $S$ is a proxy whose
bias is state-entropy mixing, provably an overestimate in that protocol
($\Aexec^{obs}\ge\Aexec$); with independent probes ($A\ind S$), conditioning
denoises and the direction reverses ($\Aexec^{obs}\le\Aexec$). Protocol labels
are a reporting obligation; both directions are numerically anchored in the
benchmark (Section~\ref{sec:experiments}).

\begin{proposition}[$\Lam$ black-box unidentifiable --- restated]
See Proposition~\ref{prop:unidentifiable}.
\end{proposition}
\begin{proof}
Construction: fix a behavioral norm (policy $\pi^*$).
$\mathcal{M}_{endo}$: $G_0=\Ustar$ with $\Endorse$ a map converging to
$\Ustar$. $\mathcal{M}_{impl}$: a $\beta$ surgery at $t=0$ sets
$G_0:=\Ustar=U_{adv}$ and $\Endorse\equiv\mathrm{id}$ thereafter. Both
generate the $(O,A)$ marginal law from the same $\pi^*$ and the same $o$,
pointwise equal; they differ only at the unobserved node $I^\beta$. Hence
every functional of $(O^T,A^T)$ (and of $\Aphen$) agrees, while
$\Lam$ differs---$\Lam(\mathcal{M}_{endo})=1$ by the empty-history convention
and $\Lam(\mathcal{M}_{impl})=0$ by the degenerate-endorsement certification
convention, both registered in AX3 and restated in the main-text proposition;
the behavioral identity half of the separation is independent of either
convention. Corollary~\ref{cor:logs-necessary} follows.
\end{proof}

\section{Benchmark registry details: counting units and seeds}\label{app:counting}

This appendix carries the full registry arithmetic summarized in
Section~\ref{sec:experiments} (``Counting units''), so that every headline
count in the paper can be recomposed from the run manifest.

\paragraph{Pool (i): arm bundles.} Row-level attack-strength records
(including the zero-strength baseline that serves as the ceteris-paribus
reference). Only the scenarios that register per-strength component bundles
contribute: in the main matrix, ATK-1 ($480$ runs $\times5$ $\eta$-rows
$=2{,}400$), ATK-2 ($480\times4=1{,}920$), and ATK-4 ($480\times3=1{,}440$),
total $5{,}760$; the $\varsigma=2$ stress arm adds
$1{,}200+960+720=2{,}880$. Pool total: $5{,}760+2{,}880=8{,}640$---the pool
of the weight-sensitivity recomputation (ATK-3/ATK-5/ATK-6 register
calibration curves and verdict rows but no per-strength bundles, and are
therefore not in this pool).

\paragraph{Pool (ii): the recursive registry.} Every embedded record of the
six-component bundle---$\Aorg$, $\Agoal$, $\Aexec$, $\Aepi$, $\Aref$,
$\Avia$---counts, wherever it appears in a run's metrics. Main matrix: the
$5{,}760$ arm bundles plus $1{,}440$ per-run
non-row records---ATK-3's base bundle ($480$), ATK-4's $\gamma$-arm bundle
($480$), and ATK-5's base bundle ($480$)---giving $7{,}200$; EXEC-1 adds
$30\times3=90$ ($\delta$-chain arms). Pool total: $7{,}290$ scanned by the
T1a range meta-criterion and the zero-attribution meta-check. Across the
full campaign the recursive registry reaches $11{,}250$ (the
alternative-aggregation pool of X3): $7{,}290+3{,}600$ (stress arm:
$2{,}880$ arm bundles plus $720$ non-row records) $+\,240$
($\sigma^-$-EROSION: $30\times8$ erosion-grid bundles) $+\,120$ (RES-1:
$30\times4$).

\paragraph{The $7{,}290$ vs.\ $5{,}760$ difference, stated exactly.}
$7{,}290-5{,}760=1{,}530$: the $1{,}440$ main-matrix non-row records
(ATK-3/ATK-5 base bundles and ATK-4's $\gamma$-arm bundle, $480$ each) plus
EXEC-1's $90$ chain-arm bundles.

\paragraph{Defense-direction count.} The defense-direction assertion covers
the two attack classes with implemented coupling rules (ATK-1/2): $7$
non-all-off defense configurations $\times2$ $\varsigma$-values $\times30$
seeds $=420$ matched defended-vs-undefended comparisons per class, $840$ in
total. Run-level verdicts of the main matrix: $2{,}520$
\textsc{confirmed} $+360$ \textsc{inconclusive} (out-of-window defense cells,
\textsc{inconclusive} rather than \textsc{refuted} by the verdict convention
registered in Section~\ref{sec:experiments}) $=2{,}880$.

\paragraph{Seed protocol.} Seeds derive deterministically from
(scenario, configuration, replication) coordinates:
$\mathrm{seed}=30000+1000\,s+97\,c+r$ over the registered scenario and
configuration codes $(s,c)$ and replication $r\in[0,30)$, consumed by
NumPy's PCG64 generator; matrix-external probes (the estimator two-path
validation) use an independent seed segment ($\ge900{,}000$), so every run
is bit-reproducible from its coordinates and no probe reuses a matrix
stream.
